\section{The group input and the algebraic data}\label{sec:reduction}

We first enlarge the input group so that each of its elements is a
commutator in the larger group. This will allow us to embed the input
group in an elementary matrix group in Section~\ref{sec:action}.
We then state the algebraic construction used in the next three
sections. All rings and algebras have identities; a homomorphism
explicitly called nonunital is not required to preserve them.
Our commutator convention is \([a,b]=aba^{-1}b^{-1}\).

\begin{lemma}\label{lem:group-input}
Let \(G\) be finitely generated with decidable word problem. There is a
finitely generated group \(\widehat G\) with decidable word problem
containing \(G\times G\), and an element \(s\in\widehat G\), such that
\[
 [(g,g),s]=(g,1)\qquad(g\in G).
\]
\end{lemma}
\begin{proof}
Take the HNN extension
\[
 \widehat G=
 \left\langle G\times G,s\ \middle|\
 s(g,g)s^{-1}=(1,g)\quad(g\in G)\right\rangle .
\]
The associated subgroups are the diagonal and the second factor, with
the indicated isomorphism. Britton's Lemma embeds the base
\cite{Britton63}; see also the precise effective HNN formulation in
\cite[p.~44]{BH74}.

We spell out the decision procedure. Store each base word as a pair of
words on the given generators of \(G\). Membership of \((a,b)\) in the
diagonal is decided by testing \(ab^{-1}=1\), and membership in the second
factor by testing \(a=1\). Replace each available pinch by
\[
 s(a,a)s^{-1}\longmapsto(1,a),
 \qquad
 s^{-1}(1,b)s\longmapsto(b,b),
\]
and multiply adjacent base words. Each pinch removes two stable letters.
When no pinch remains, a word still containing a stable letter is
nontrivial by Britton's Lemma. Otherwise decide the resulting base word
coordinatewise. Thus the procedure terminates and decides the word
problem, without requiring a finite presentation of \(G\).
The two copies of a finite generating set, together with \(s\), generate
\(\widehat G\). Finally,
\[
 [(g,g),s]=(g,g)\,s(g^{-1},g^{-1})s^{-1}
           =(g,g)(1,g^{-1})=(g,1).
\]
This includes finite groups and the trivial group; in the latter case
\(\widehat G\) is infinite cyclic.
\end{proof}

Here is the algebraic construction. A simple algebra means a nonzero
algebra with no proper nonzero two-sided ideal. A subalgebra with identity
\(e\) need not contain the identity of the ambient algebra; its units
are taken relative to \(e\). For an idempotent \(q\in B\), the
subalgebra \(qBq=\{qbq:b\in B\}\) is called a \emph{corner};
its identity is \(q\).

\begin{theorem}\label{thm:algebra-package}
For every finitely generated group \(L\) with decidable word problem,
there are the following data over \(k=\F_2\):
\begin{enumerate}[label=\textup{(\roman*)}]
\item A nonzero finitely presented unital algebra \(B\), a simple
subalgebra \(A\subseteq B\) with identity \(e\), and an embedding
\(L\hookrightarrow A^\times\).
\item An injective nonunital algebra homomorphism \(\phi:B\to B\) with
\(\phi(B)\subseteq A\) and \({p=\phi(1_B)\ne1_B}\).
\item Elements \(s_0,s_1,t_0,t_1\in pBp\) commuting with \(\phi(B)\)
and satisfying
\[
 t_i s_j=\delta_{ij}p,\qquad s_0t_0+s_1t_1=p
 \quad(i,j\in\{0,1\}).
\]
\item An element \(b_*\in B\) with the following property. Put
\[
 R=B\otimes_k B^{\op},\qquad
 \psi=\phi\otimes\phi^{\op},
 \qquad (a\otimes b^{\op})\cdot m=amb.
\]
Then
\[
 r\ne0\ \Longrightarrow\ \psi(r)\cdot b_*\ne0
 \qquad(r\in R).
\]
\end{enumerate}
\end{theorem}

Theorem~\ref{thm:algebra-package} is proved in
Section~\ref{sec:algebra}, using the monoid construction of
Section~\ref{sec:bootstrap} and Appendix~\ref{sec:compiler}.
To see how these data produce an action, regard the elements of \(R\)
as coefficients for row operations on \(B^6\). Such a coefficient
acts on a coordinate by a finite sum of left and right multiplications.
Simplicity of \(A\) will let a nonzero coordinate produce any prescribed
element of \(A\); applying \(\phi\) places all coordinates in that
subalgebra. Finite presentation of \(B\) supplies a finitely presented
Steinberg group acting by row operations. The splitting pairs make its matrix
kernel disappear under the induced coefficient map, while \(b_*\)
detects every remaining nonidentity matrix after that map is applied.
Finally, the support identity \(p\ne1_B\) makes successive vector
stages strictly larger, so that shifting stages gives a nontrivial
permutation. Sections~\ref{sec:affine}--\ref{sec:action} establish these
claims and assemble the action.

In particular, \(p\phi(b)=\phi(b)p=\phi(b)\), so
\(1_B\notin\phi(B)\). The detection condition concerns the action
after applying \(\psi\); it does not require the original action of
\(R\) on \(B\) to be faithful.

Apply the theorem with \(L=\widehat G\) from
Lemma~\ref{lem:group-input}. For the remainder of
Sections~\ref{sec:affine}--\ref{sec:action}, fix these data. Tensor
products are over \(k\). Besides \(p\), define
\[
 P=p\otimes p^{\op},\qquad
 U_i=s_i\otimes p^{\op},\qquad
 V_i=t_i\otimes p^{\op}.
\]
Tensoring injective linear maps over a field shows that \(\psi\) is
injective, so \(P\ne0\). Also \(P\ne1_R\), since the nonzero tensor
\((1_B-p)\otimes1_B\) annihilates \(P\) on the left.
The elements \(U_i,V_i\) lie in \(PRP\), commute with
\(\psi(R)\), and satisfy
\begin{equation}\label{eq:ring-pairs}
 V_iU_j=\delta_{ij}P,\qquad U_0V_0+U_1V_1=P.
\end{equation}
These identities follow directly from the corresponding identities
in \(B\); the second tensor factor is always the idempotent \(p^{\op}\).

The final passage from an action to a simple group uses the following
established result, in the precise form needed here.

\begin{theorem}[Zaremsky]\label{thm:zaremsky}
Suppose a finitely presented group \(\Lambda\) acts faithfully on a
nonempty countable set \(X\). If the point stabilizers are finitely
generated and there are finitely many orbits on the unordered
two-element subsets of \(X\), then \(\Lambda\) embeds in a nontrivial
finitely presented simple group.
\end{theorem}

These are the hypotheses of a type-\textup{(A)} action in
\cite[Theorem~A and Corollary~B]{Zaremsky2024Finite}.
Finite presentation follows from that theorem; the canonical embedding
and simplicity are those of the twisted Brin--Thompson construction
\cite[Section~3 and Theorem~3.4]{BelkZaremsky2022}. Since \(X\) is nonempty, the
constructed group contains Thompson's group \(V\) on a coordinate and
is nontrivial. Our construction will provide a countably infinite \(X\)
with one orbit on pairs, and will prove each of the other hypotheses
separately.
